\section{模型设计：统一感知、规划与执行的多阶段架构}

\subsection{现有 GUI Agent 的关键短板}
讲者归纳了当前系统的三类问题：视觉 grounding 不稳、长程规划弱、跨平台泛化差。很多模型在短任务上可用，但在高复杂度任务中会快速退化。

\begin{importantbox}{三类能力缺口}
\begin{enumerate}
    \item \textbf{Grounding}：看到了但定位不准，或目标解析歧义；
    \item \textbf{Planning}：会执行单步，但缺少全局子目标管理；
    \item \textbf{Generalization}：换界面、换任务模板后性能骤降。
\end{enumerate}
\end{importantbox}

\subsection{课程提出的整合思路}
从字幕可见，讲者强调通过统一 AR 风格的生成范式，将视觉理解、推理链（monologue/reasoning trace）和动作生成纳入同一训练接口。这样可减少多模型拼接带来的接口损耗。

\begin{figure}[H]
\centering
\includegraphics[width=0.88\textwidth]{frame-07-model-architecture.jpg}
\caption{视频约 01:16:05：模型架构强调统一输入输出格式与多阶段训练目标}
\end{figure}

\begin{knowledgebox}{多阶段训练的工程收益}
\begin{itemize}
    \item 阶段 1：强化视觉 grounding 与界面语义理解；
    \item 阶段 2：引入推理链与规划信号，提升长程决策；
    \item 阶段 3：对齐执行策略与奖励反馈，优化任务完成率。
\end{itemize}
\end{knowledgebox}

\subsection{Planner 信号的引入}
课程后段反复提到 planner 对难任务收益显著。直观上，planner 不是替代动作模型，而是提供可解释的中间结构，让模型在长任务中减少盲目试探。

\begin{warningbox}{Planner 不等于 ``额外冗余''}
若任务本身是长链路且状态复杂，缺 planner 往往意味着更多无效动作与回退。额外规划开销通常小于错误执行成本。是否使用 planner 应按任务复杂度分层决策。
\end{warningbox}

\subsection{本章小结}
模型架构的关键不在参数规模，而在 ``统一接口 + 分阶段训练 + 规划信号''。这三点共同决定了系统在高复杂度任务中的稳定性。

